\section{Limitations and Broader Considerations}

Our theoretical results are local and conditional. Ordering by
success-per-cost is optimal only under the stated sequential model; real proof
actions have dependent outcomes, can reveal information, and can change the
future action space. The controller's richer score is therefore a policy class
to evaluate, not a proof of global optimality.

Verifier-derived progress is more objective than unconstrained model
self-evaluation but is not synonymous with distance to a proof. A candidate
with fewer goals may have replaced an easy goal with a much harder one, and a
decomposition may create formally valid but strategically useless helpers.
Calibration must be checked across task domains and model families.

Structured search has overhead. It uses longer prompts, executes control logic,
and may spend early budget constructing obligations that a direct model could
solve immediately. This is why minimal refinement is retained as a first-class
baseline and why we report task-stratified crossover behavior rather than
claiming that structured mode should be universal.

The current system targets Lean 4. The proof-system adapter is narrow enough to
make another backend plausible, but obligation semantics, diagnostics, tactic
granularity, and checker cost differ across assistants. Results should not be
generalized to Isabelle, Rocq, or Metamath without new experiments.

Finally, a checker-accepted proof certifies the formal statement, not the
faithfulness of an automatically formalized statement to the user's intent.
The primary experiments fix formal statements, and the end-to-end setting
reports formalization separately. Statement preservation and anti-cheating
checks prevent several mechanical shortcuts, but they do not replace human
review of high-stakes specifications.
